\section{Multiplication operators in characteristic two}
\label{sec:modular-multiplication}

The failure of equivariant averaging in characteristic two leaves room
for constructions that use nilpotent operators. We describe one such
space explicitly. It contains finite decompositions whose cost lies
below the square-root-of-two benchmark, but its commuting pair
subalgebra has a uniform obstruction to asymptotic improvement.
Throughout this section the field \(K\) has characteristic two.

\begin{proposition}[Multiplication and partial transpose]
\label{prop:modular-multiplication}
Put
\[
 A_k=K[x_1,y_1,\ldots,x_k,y_k]/(x_i^2,y_i^2),
 \qquad U=(K^2)^{\otimes k},\qquad q=2^k.
\]
Let \(E|1\rangle=|0\rangle\), \(E|0\rangle=0\), and \(F=E^{\mathsf T}\).
Identify \(A_k\) with \(U\otimes U\) by making
\[
 x_i=E\otimes I+I\otimes E,\qquad
 y_i=F\otimes I+I\otimes F
\]
on the \(i\)-th pair and taking the product of the cyclic vectors
\(|01\rangle\). Regroup the first and second spinors into the two
copies of \(U\). Write \(m_f\) for multiplication by \(f\), and set
\[
 J=\begin{pmatrix}0&1\\1&0\end{pmatrix}^{\otimes k},
 \qquad
 \Gamma_J=\mathrm{id}\otimes
       \bigl(T\mapsto J^{-1}T^{\mathsf T}J\bigr).
\]
Then
\begin{equation}\label{eq:multiplication-pt}
 \Gamma_J(m_f)=m_{\mathcal Lf},
 \qquad
 \mathcal L=\prod_{i=1}^k
           (1+\partial_{x_i}\partial_{y_i}).
\end{equation}
Consequently \(f+\mathcal Lg=1\) gives
\(I_{q^2}=m_f+\Gamma_J(m_g)\). Conjugating \(m_g\) by \(I\otimes J\)
converts this to ordinary partial transpose and preserves both ranks.
\end{proposition}
\begin{proof}
The displayed \(x,y\) commute and square to zero. On one pair the
images of \(1,x,y,xy\) applied to the cyclic vector are
\[
 |01\rangle,\quad |00\rangle,\quad |11\rangle,\quad
 |01\rangle+|10\rangle.
\]
They form a basis, proving the regular-module identification.
On the corresponding four multiplication operators, \(\Gamma_J\)
fixes \(1,x,y\) and sends \(xy\) to \(xy+1\). This is exactly
\(1+\partial_x\partial_y\). Tensoring the local identities proves
\eqref{eq:multiplication-pt}. Finally \(J=J^{-1}=J^{\mathsf T}\),
which gives the stated conversion to ordinary partial transpose.
\end{proof}

Put \(p_i=x_i y_i\). In the spinor coordinates,
\(m_{p_i}=I+F_i\), where \(F_i\) exchanges the two spinors at site \(i\).
For a polynomial in the \(p_i\), the map \(\mathcal L\) is the
linear substitution
\[
 \tau(p_S)=\prod_{i\in S}(1+p_i).
\]
Substitution here is applied to the unique multilinear
representative; it is not a homomorphism of the square-zero algebra.

\begin{corollary}[Two finite decompositions]
\label{cor:char2-eight-wire}
\label{cor:char2-thirty-two-wire}
Let \(s_j\) denote the elementary symmetric polynomial of degree \(j\)
in the indicated pair variables. For \(k=3\), take \(f=s_2\);
for \(k=5\), take \(f=s_3+s_4\). In both cases
\(\mathcal Lf+f=1\). The multiplication ranks are respectively
\(8\) and \(66\). Thus there are ordinary partial-transpose
decompositions of \(I_{64}\) and \(I_{1024}\) with rank pairs
\[
 (8,8)\qquad\hbox{and}\qquad(66,66).
\]
Their total costs \(16\) and \(132\) are below
\(8^{\sqrt2}\) and \(32^{\sqrt2}\), respectively.
These are finite statements.
\end{corollary}
\begin{proof}
The binomial formula for \(\tau(s_j)\) gives
\[
 \tau(s_2)=1+s_2\quad(k=3),\qquad
 \tau(s_3)=s_2+s_3,\quad
 \tau(s_4)=1+s_2+s_4\quad(k=5).
\]
For the ranks, decompose each local algebra as a module for \(p_i\):
\[
 K[1,p_i]\ \oplus\ Kx_i\ \oplus\ Ky_i.
\]
A set of \(r\) active pairs therefore contributes \(2^{k-r}\)
copies of the regular square-zero algebra on those \(r\) pair variables.
For \(k=3\), the two-active-pair sectors contribute \(6\), and the
three-active-pair sector contributes \(2\). The latter multiplication
has rank one from degree zero to degree two and rank one from degree
one to degree three.

For \(k=5\), sectors with fewer than three active pairs contribute zero.
The three-active-pair sectors contribute
\(\binom53 2^2=40\). On four active pairs, multiplication by
\(s_3+s_4\) has rank two, contributing \(\binom54 2\cdot2=20\).
On five active pairs, the degree-zero image has a nonzero degree-three
part. The five degree-one images have degree-four coefficient matrix
\(I_5+\mathbf1\mathbf1^{\mathsf T}\), of rank four. Its
one-dimensional kernel is generated by \(\mathbf1\), and that
combination of the degree-one images is the nonzero top monomial.
Thus these images have rank five. All degree-two images lie in the
same top-monomial line, and higher input degrees vanish. The full
five-pair sector has rank six, giving \(40+20+6=66\).
All ranks in this argument are unchanged on extending \(K\).
The benchmark comparisons are exact: \(4<3\sqrt2\), while
\(24^2<2\cdot17^2\) and \(132^{17}<2^{120}\) imply
\(132<2^{120/17}<32^{\sqrt2}\).
\end{proof}


The inclusion-matrix rank used below is a special case of Wilson's
integral diagonal form~\cite{Wilson1990}; a precise statement and proof
appear in Ghorbani, Khosrovshahi, Maysoori and Mohammad-Noori,
Corollary~3, printed page~10~\cite{Ghorbani2007}.
For an inclusion matrix from \(r\)-subsets to \((r+d)\)-subsets of an
\(m\)-set with \(2r+d\le m\), its diagonal entries are
\(\binom{r+d-i}{r-i}\), with multiplicity
\(\binom mi-\binom m{i-1}\), for \(0\le i\le r\).
When \(d\) is a power of two and \(r<d\), all these entries are odd.
We give an elementary proof of this special case as part of the rank
calculation, so no general diagonal-form theorem is needed below.

\begin{proposition}[Exact ranks of a binary-degree family]
\label{prop:elementary-pair-rank}
Let \(d\) be a power of two, \(k=2d-1\), and \(f=s_d(p_1,\ldots,p_k)\).
Then \(f+\mathcal Lf=1\), and the rank of multiplication by \(f\) is
\begin{equation}\label{eq:elementary-pair-rank}
 R_{k,d}=\binom{2k}{k-d}
          +2\sum_{j=0}^{k-d-1}\binom{2k}{j},
\end{equation}
where an empty sum is zero. Along this sequence,
\[
 \frac{\log_2(2R_{k,d})}{k}
 \longrightarrow 2H_2(1/4)
 =4-\frac32\log_2 3>\frac32>\sqrt2.
\]
Here \(H_2(u)=-u\log_2u-(1-u)\log_2(1-u)\).
\end{proposition}
\begin{proof}
The coefficient of \(s_i\) in \(\tau(s_d)\) is
\(\binom{2d-1-i}{d-i}\). It is one for \(i=0,d\) and zero for
\(0<i<d\). Indeed, in characteristic two,
\[
 (1+z)^{2d-1-i}=(1+z^d)(1+z)^{d-1-i},
\]
whose coefficient of \(z^{d-i}\) is zero for \(0<i<d\).
This proves the self equation.

On \(m\) active pair variables, multiplication by \(s_d\) maps
degree \(r\) to degree \(r+d\) by the subset-inclusion matrix.
Suppose first \(2r+d\leq m\). Restrict to any universe of
\(2r+d\) coordinates and index target rows by their complementary
\(r\)-subsets. The resulting square matrix \(D\) has entry one
exactly on disjoint \(r\)-subsets. Thus
\[
 (D^2)_{S,T}=\binom{d+|S\cap T|}{r}.
\]
Since \(r\leq(m-d)/2<d\), Frobenius gives
\(\binom{d+u}{r}=\binom ur\) in \(K\). Consequently \(D^2=I\).
A kernel vector for the larger inclusion matrix restricts to a
kernel vector on every such universe, so it vanishes.
Complementing subsets proves the dual full-row-rank case.
Each degree block therefore has rank
\(\binom m{\min(r,m-d-r)}\).

The active-sector decomposition gives
\[
 R_{k,d}=\sum_{m=d}^k\binom km2^{k-m}
       \sum_{r=0}^{m-d}\binom m{\min(r,m-d-r)}.
\]
The inner sum is the coefficient of \(z^d\), plus twice the sum
of coefficients of \(z^\ell\) for \(\ell>d\), in
\((z+z^{-1})^m\). Summing the active-sector weights replaces this
Laurent polynomial by
\[
 (z+2+z^{-1})^k=z^{-k}(1+z)^{2k}.
\]
Its coefficient of \(z^\ell\) is \(\binom{2k}{k-\ell}\),
proving \eqref{eq:elementary-pair-rank}. The integer coefficients
in this counting step are ordinary dimensions, not field elements.

The binomial tail differs from its largest term by at most a factor
\(2k+1\). Stirling's formula and
\((k-d)/(2k)\to1/4\) give the stated limit.
Finally \(3^3<2^5\) implies
\(4-\frac32\log_2 3>\frac32\).
\end{proof}


\begin{proposition}[Forms supported above the middle degree]
\label{prop:high-degree-multiplication}
Restrict \(f\in A_k\) to monomials of ordinary total degree greater
than \(k\). Then \(f+\mathcal Lf=1\) has no solution for even \(k\).
For odd \(k\) it has exactly one solution. This solution belongs to
the pair algebra and its Boolean evaluation is
\[
 f(1_S)=[\,|S|>k/2\,].
\]
No symmetry assumption on \(f\) is required.
\end{proposition}
\begin{proof}
Decompose the monomial basis by the sites occupied by exactly one
of \(x_i,y_i\). Fix such a choice on \(u\) sites. The remaining
\(m=k-u\) sites carry a polynomial \(h\) in their pair variables,
and the degree restriction says that every term of \(h\) has pair
degree greater than \(m/2\). The operator \(\mathcal L\) preserves
this sector and acts on it by \(\tau\).

If \(u>0\), the equation in this sector is \(h+\tau h=0\).
The degree restriction makes \(h(1_S)=0\) for \(|S|\leq m/2\).
Complement invariance makes it zero on the other Boolean points
as well. Boolean evaluation determines a multilinear polynomial,
so \(h=0\). Thus every solution lies in the pair sector \(u=0\).

In that sector, evaluation gives
\(f(1_S)+f(1_{S^{\mathrm c}})=1\).
For even \(k\), both values vanish on the middle layer, a
contradiction. For odd \(k\), low-weight values vanish and the
complement equation forces all high-weight values to be one.
There is a unique multilinear polynomial with these values.
Its coefficients in degrees at most \(\lfloor k/2\rfloor\) vanish,
by triangular Boolean interpolation from its zero low-weight
values. Its ordinary degree is therefore greater than \(k\),
and its complementary values prove \(f+\tau f=1\).
\end{proof}

A finite cost below \(q^{\sqrt2}\) does not by itself supply an improving
sequence: tensoring a two-cut decomposition creates mixed cut patterns.
The following theorem controls all tensor powers within the pair
algebra, including arbitrary coefficients that do not factor between
copies.

\begin{theorem}[Projected pair-polynomial rank bound]
\label{thm:projected-pair-barrier}
Let \(S\subseteq\{0,1\}^k\) be nonempty and stable under bitwise
complement, put \(Q=|S|\), and let \(U_S\) be its coordinate span.
Let \(P\) project onto \(U_S\otimes U_S\), and write
\[
 M_f=P f(p_1,\ldots,p_k)P.
\]
Let \(J\) be the complement permutation on \(U_S\).
Suppose only that
\[
 I_{U_S\otimes U_S}=M_f+\Gamma_J(M_g).
\]
For \(A\subseteq[k]\), define
\[
 w(A)=\bigl|\{a|_{A^{\mathrm c}}:
          a\in S,\ a\mathbin{\triangle}A\in S\}\bigr|.
\]
Here binary strings are identified with subsets. Then
\begin{equation}\label{eq:projected-pair-finite}
 \operatorname{rank}M_f+\operatorname{rank}M_g
 \geq \sum_{A\subseteq[k]}\min\{w(A),w(A^{\mathrm c})\}.
\end{equation}
On \(t\) copies of \(U_S\), allow arbitrary polynomials in all \(kt\)
pair variables. Every corresponding projected identity has cost at least
\begin{equation}\label{eq:projected-pair-words}
 B_{S,t}=
 \sum_{A_1,\ldots,A_t}
 \min\left\{\prod_iw(A_i),\prod_iw(A_i^{\mathrm c})\right\}.
\end{equation}
For each fixed \(S\),
\begin{equation}\label{eq:projected-pair-rate}
 \lim_{t\to\infty}B_{S,t}^{1/t}
 =D_S:=\sum_A\sqrt{w(A)w(A^{\mathrm c})}
 \geq Q^{3/2}.
\end{equation}
\end{theorem}
\begin{proof}
The local operator \(p_i=I+F_i\) annihilates \(|00\rangle,|11\rangle\)
and preserves their complementary two-dimensional space.
Consequently pair polynomials preserve the xor \(A=a\triangle b\)
and the common coordinates \(a|_{A^{\mathrm c}}\). After projection
these index \(w(A)\) disjoint nonempty blocks. Every diagonal entry
in such a block is \(f(1_A)\), so
\[
 \operatorname{rank}M_f
 \geq\sum_Aw(A)[f(1_A)\ne0].
\]
The diagonal of \(\Gamma_J(M_g)\) at \(|a,b\rangle\) is the diagonal
of \(M_g\) at \(|a,b^{\mathrm c}\rangle\). Therefore the projected
identity gives
\[
 f(1_A)+g(1_{A^{\mathrm c}})=1
\]
whenever \(w(A)>0\). Complement stability implies
\(w(A)>0\) if and only if \(w(A^{\mathrm c})>0\).
Apply the same rank bound to \(g\), reindex by complement, and take
the smaller coefficient term by term. This proves
\eqref{eq:projected-pair-finite}. The coordinate set on \(t\) copies is
\(S^t\), whose weights factor, proving
\eqref{eq:projected-pair-words} without any factorization assumption
on the polynomials.

Discard zero terms and give the remaining subsets the probability law
\[
 \pi_A=\frac{\sqrt{w(A)w(A^{\mathrm c})}}{D_S},
 \qquad Z_A=\log\frac{w(A)}{w(A^{\mathrm c})}.
\]
Complementation preserves \(\pi\) and negates \(Z\). Thus \(Z\)
has mean zero and finite variance. For independent samples,
\[
 B_{S,t}=D_S^t\,
   \mathbb E\exp\left(-\frac12|Z_1+\cdots+Z_t|\right).
\]
The expectation is at most one. If the variance is positive,
Chebyshev's inequality bounds it below by
\(\frac12\exp(-\sqrt{2t\operatorname{Var}Z}/2)\);
if the variance is zero it equals one. This proves the limit.

Put \(T_A=S\cap(S\triangle A)\) and \(m(A)=|T_A|\).
Complement stability gives \(T_A=T_{A^{\mathrm c}}\).
The map \(a\mapsto(a|_A,a|_{A^{\mathrm c}})\) is injective on \(T_A\),
and its two coordinates range over the patterns counted by
\(w(A^{\mathrm c})\) and \(w(A)\). Hence
\[
 m(A)\leq w(A)w(A^{\mathrm c}),\qquad
 m(A)\leq Q,\qquad \sum_A m(A)=Q^2.
\]
The last equality counts ordered pairs in \(S\). It follows that
\[
 D_S\geq\sum_A\sqrt{m(A)}
       \geq\frac1{\sqrt Q}\sum_A m(A)=Q^{3/2}.
\]
\end{proof}

The limit in \eqref{eq:projected-pair-rate} is for a fixed coordinate
set \(S\). It does not assert a uniform finite bound for arbitrary
varying coordinate sets.

\begin{corollary}[The full pair algebra]
\label{thm:pair-product-barrier}
On all \(k\) spinor pairs, every identity \(f+\tau g=1\) has cost at least
\[
 \sum_{j=0}^k\binom kj2^{\min(j,k-j)}.
\]
This quantity lies between
\(2^{k+\lfloor k/2\rfloor}/(k+1)\) and \(2^{3k/2}\).
Thus the pair algebra has lower exponent \(3/2\) uniformly as \(k\)
varies, over every characteristic-two extension field.
\end{corollary}
\begin{proof}
For the full cube, \(w(A)=2^{k-|A|}\), so the first claim follows
from \eqref{eq:projected-pair-finite}. The middle binomial coefficient
is at least \(2^k/(k+1)\), giving the lower bound; bounding every
factor \(2^{\min(j,k-j)}\) by \(2^{k/2}\) gives the upper bound.
\end{proof}

\begin{proposition}[The balanced compound is a projected pair family]
\label{prop:balanced-pair-identification}
Let \(T(z)\) be the scalar balanced compound of
Section~\ref{sec:balanced-compound}, with \(2n\) sites and
weight-\(n\) endpoints. On the same endpoint spaces put
\(F(t)=\prod_{i=1}^{2n}(I+t p_i)\). Then
\begin{equation}\label{eq:balanced-pair-identification}
 F(t)_{A,B,C,D}=T(1+t^2)_{A,B^{\mathrm c},C^{\mathrm c},D}.
\end{equation}
Consequently their coefficient spans agree after these fixed endpoint
permutations. The odd coefficients of \(F(t)\) vanish on this projection.
\end{proposition}
\begin{proof}
After complementing the second and third binary indices, the
one-site minor tensor has the following nonzero entries:
\[
\begin{array}{c|rrrrrr}
 (a,b,c,d)&0000&1111&0101&1010&0110&1001\\ \hline
 T(z)&1&1&1&z&1&1+z.
\end{array}
\]
Set \(z=(1+t)^2\). Multiplying an index-one coordinate on the four
legs by \(1,t(1+t),(1+t)^{-1},t^{-1}\), respectively, changes these
entries to \(1,1,1+t,1+t,t,t\), exactly those of \(I+t p_i\).
This calculation is over \(K(t)\). On four weight-\(n\) endpoints
the total multiplier is
\[
 [1\cdot t(1+t)\cdot(1+t)^{-1}\cdot t^{-1}]^n=1.
\]
The multipliers therefore disappear exactly, proving the polynomial
identity \eqref{eq:balanced-pair-identification}.
Substitution \(z=1+t^2\) is triangular and invertible on the two
coefficient spaces.

For completeness, these permutations also transport the cut
convention. With \(R=I\otimes J\), \(C=J\otimes I\), and \(M=RTC\),
ordinary second-factor transpose satisfies
\[
 \Gamma_2(M)=\Gamma_2(T)(J\otimes J),
 \qquad \Gamma_J(M)=R\Gamma_2(T)C.
\]
Thus both ranks are preserved and the crossing at \(z=1\) becomes
the identity at \(t=0\).
\end{proof}

The general bound therefore contains the scalar-compound
tensor-power obstruction as a special case. It also allows arbitrary
nonsymmetric pair-polynomial coefficients on those projected spaces.
It does not cover the full algebra in the separate \(x_i,y_i\),
all diagram operators, or unrestricted crossing decompositions.

\subsection{Two triangular graph ideals}

Another way to impose low multiplication rank is to take a principal
ideal obtained by a triangular change of variables. The following
obstruction allows two different such ideals in a common set of coordinates.

\begin{proposition}\label{prop:triangular-graph-ideals}
Let \(K\) have characteristic two, put
\[
 B=K[y_1,\ldots,y_k]/(y_i^2),\qquad
 A=B[x_1,\ldots,x_k]/(x_i^2),\qquad
 L=\prod_i(1+\partial_{x_i}\partial_{y_i}),
\]
and let \(F_i,G_i\in B\) have zero constant term. For arbitrary
\(u,v\in B\), set
\[
 f=u\prod_i(x_i+F_i),\qquad g=v\prod_i(x_i+G_i).
\]
Then \(\operatorname{rank}_A(m_f)=\operatorname{rank}_B(m_u)\leq2^k\).
If \(k\geq4\), no such pair satisfies \(f+Lg=1\).
\end{proposition}
\begin{proof}
Every zero-constant element of \(B\) squares to zero. Thus the
substitution \(x_i\mapsto x_i+F_i\), fixing \(B\), is an involutive
algebra automorphism. It reduces the rank assertion to multiplication
by \(u x_1\cdots x_k\), whose image is a copy of \(uB\).

Suppose that \(f+Lg=1\). Its coefficient of \(x_1\cdots x_k\)
gives \(u=v\). First suppose that this common element is a nonunit.
Both \(f\) and \(g\) then have minimum ordinary degree greater than
\(k\). The proof of Proposition~\ref{prop:high-degree-multiplication}
shows that the only \(L\)-fixed polynomial of minimum degree greater
than \(k\) is zero: in each singleton-occupancy sector its Boolean
values vanish below the middle layer, and complement invariance kills
the remaining values. Applying \(L\) to the equation shows that
\(f+g\) is fixed, hence \(f=g\). The coefficient of
\(\prod_{j\ne i}x_j\) in \(f+Lf=1\) is \(\partial_{y_i}u\).
Consequently \(u\) is constant and, being a nonunit, zero, a contradiction.

It remains to treat unit \(u\). Write \(X=\prod_i x_i\) and
\(X_i=\prod_{j\ne i}x_j\). Complementation of the squarefree
\(x\)-indices conjugates \(L\) to the automorphism
\(\sigma:y_i\mapsto y_i+x_i\), fixing the \(x_i\). The equation becomes
\[
 H+\sigma K=X,\qquad
 H=u\prod_i(1+x_iF_i),\qquad K=u\prod_i(1+x_iG_i).
\]
Both \(H,K\) are units. Direct differentiation, using \(F_i^2=0\), gives
\[
 \frac{\partial_{x_i}(H+X)}{H+X}=F_i+u^{-1}X_i.
\]
For example, \((H+X)^{-1}=H^{-1}+XH^{-2}\), and the two terms
containing \(XF_i\) cancel. The chain rule for \(\sigma\), followed
by another application of \(\sigma\), therefore yields
\begin{equation}\label{eq:graph-ideal-logarithm}
 G_i+u^{-1}\partial_{y_i}u
       +\sum_jx_j\partial_{y_i}G_j
 =F_i(y+x)+u(y+x)^{-1}X_i.
\end{equation}
The left side has \(x\)-degree at most one. Since \(k\geq4\), the
last term on the right has \(x\)-degree at least three. Comparing
each coefficient \(x_ax_b\), \(a\ne b\), gives
\(\partial_{y_a}\partial_{y_b}F_i=0\). Thus every \(F_i\) is linear.
Now the coefficient of \(X_i\) in \eqref{eq:graph-ideal-logarithm}
is zero on the left and \(u^{-1}\ne0\) on the right, the final
contradiction.
\end{proof}

The cutoff is sharp for this family: at \(k=1,2,3\), respectively,
the self pairs with
\[
 f=(1+y_1)x_1,\qquad
 f=x_2(x_1+y_2),\qquad
 f=\prod_{i=1}^3\left(x_i+\prod_{j\ne i}y_j\right)
\]
have \(f+Lf=1\) and ranks \(2,4,8\). The third expression expands to
\(x_1x_2x_3+s_2(p_1,p_2,p_3)\).
The proposition excludes these triangular factors, not arbitrary
low-rank elements of \(A\) or unrestricted partial-transpose identities.

